\documentclass[10pt,a4paper]{amsart}
\usepackage[margin=19mm]{geometry}
\usepackage{amsmath,amssymb,mathtools}
\usepackage[scr=rsfs,cal=euler]{mathalfa}
\usepackage{xfrac}
\usepackage[unicode,hidelinks,bookmarks=false]{hyperref}
\newcommand{\ie}{i.e.\ }
\newcommand{\cf}{cf.\ }
\newcommand{\resp}{respectively\ }
\newcommand{\Subsection}{\subsection}
\providecommand{\nobreakdash}{\nobreak\hbox{-}}
\setlength{\emergencystretch}{2em}
\multlinegap0pt
\newtheorem{lemma}{Lemma}
\usepackage{cleveref}
\crefname{lemma}{Lemma}{Lemmas}
\title{Asymptotic Transfer in Critical Recursive Composition Schemes}
\author{Michael Drmota, Z\'ephyr Salvy}
\date{Source: arXiv:2603.06185v1 (CC BY 4.0)}
\begin{document}
\maketitle

\noindent\emph{Source and licence.} Asymptotic Transfer in Critical Recursive Composition Schemes. Version arXiv:2603.06185v1 (CC BY 4.0), distributed by arXiv under the Creative Commons Attribution 4.0 International licence, http://creativecommons.org/licenses/by/4.0/, as recorded in the arXiv licence metadata for this exact version and shown on the abstract page https://arxiv.org/abs/2603.06185v1. Full licence notice: \texttt{licenses/arxiv-2603.06185-CC-BY-4.0.txt}.\par\smallskip

\noindent\textit{Editorial scope.} Counted unit: the article's lemma LeSing1 (source0.tex lines 491--503). The article's complete original proof of it is delivered with it (source0.tex lines 505--598, one \texttt{proof} environment of 94 lines). No mathematics is written, repaired or re-derived: the statement and the proof are the article's own lines, in the article's own notation, and no retained line is substituted or re-flowed.  Retained uncounted as the source-local context the proof needs: lines 178--180.  Nothing was omitted from any retained span, so the package carries no omission record.  No retained line contains a citation, so the wrapper carries no bibliography and no bibliography entry.  No retained line contains a figure environment, an image-inclusion command or drawing code, so no image is delivered and the package declares no asset dependency.  Editorial formatting only, in addition: an AMS \texttt{amsart} preamble that restates the article's own declarations and macros used by the retained lines, namely the environments \texttt{lemma} on separate counters, together with the standard packages those lines need.  The retained environments print in this document as Lemma 1 in order of appearance, whereas the article prints the same items with the numbering of its own full document.  The complete native source of the article as distributed in the arXiv e-print for this exact version is bundled unchanged as \texttt{sources/195846/source0.tex}.
\begin{equation}\label{eq32sing}
    f(z) = g(z) + h(z) (z-z_0)^{3/2},
\end{equation}

\begin{lemma}\label{LeSing1}
    Suppose $f(z)$ has a $3/2$-singularity at $z_0$ of the form 
    \cref{eq32sing} with $g'(z_0)\ne 0$ and $h(z_0) \ne 0$.
    Then $f(z)=u$ can be locally
    inverted and the inverse function has 
    a $3/2$-singularity at $u_0 = g(z_0)$:
    \begin{equation}\label{eqLeSing1}
    z = G(u) + H(u) (u-u_0)^{3/2}
    \end{equation}
    with analytic functions $G(u)$, $H(u)$ that satisfy
    $G(u_0) = z_0$, $G'(u_0) = 1/g'(z_0)$, and 
    $H(u_0) = - h(u_0)/(g'(z_0))^{5/2}$.
\end{lemma}

\begin{proof}
It is very easy to check this property in a formal way. By setting
$Z = z-z_0$ the relation \cref{eq32sing} rewrites (with $u = f(z)$) to
\[
u = g_0 + g_2 Z + g_3 Z^{3/2} + g_4 Z^2 + \cdots,
\]
where $g_0 = g(z_0)$, $g_2 = g'(z_0) \ne 0$, and $g_3 = h(z_0) \ne 0$.
Setting $U = u-g_0$ we also have
\begin{equation}\label{eq32singvar}
U = g_2 Z + g_3 Z^{3/2} + g_4 Z^2 + \cdots
\end{equation}

If we now assume that the inverted function expands in a corresponding
way, then we should have a relation of the form
\[
Z = h_2 U + h_3 U^{3/2} + h_4 U^2 + \cdots
\]
with $h_2\ne 0$ and $h_3\ne 0$.
From that it would follow that
\[
Z^{3/2} = h_2^{3/2} U^{3/2}\left( 1 + \frac{3h_3}{2h_2} U^{1/2} + \cdots \right).
\]
By inserting the representation for $Z$ and $Z^{3/2}$ into \cref{eq32singvar}
and by comparing the coefficients of $U$ and $U^{3/2}$ it follows that
$
g_2 h_2 = 1 \quad \mbox{and} \quad g_2 h_3 + g_3 h_2^{3/2}
$
which implies that
$
h_2 = \frac 1{g_2} \quad \mbox{and} \quad h_3 = -{g_3}/{g_2^{5/2}}.
$

The main difficulty in the proof is to show that there is representation
of the form \cref{eqLeSing1} for the inverse function.
    We set $f(z) = u$ and rewrite \cref{eq32sing} to
    \[
    Q(z,u) = (u- g(z))^2 - h(z)^2 (z-z_0)^3 = 0.
    \]
    At $z=z_0$ and $u = u_0 = g(z_0)$ we have
    \[
    Q(z_0,u_0) = Q_{z}(z_0,u_0) = 0, \quad  Q_{zz}(z_0,u_0) = 2g'(z_0)^2 \ne 0.
    \]
    Hence, by the Weierstrass theorem we can represent $Q(z,u)$ locally as
    \[
    Q(z,u) = K(z,u)( (z-z_0)^2 + h_1(u)(z-z_0) + h_0(u)),
    \]
    where $h_0(u), h_1(u), K(z,u)$ are analytic and satisfy
    $h_0(u_0) = h_1(u_0) = 0$ and $K(z_0,u_0) \ne 0$.
    Hence, if we want to invert the function $u = f(z)$
    we have to solve the quadratic equation
    \[
    (z-z_0)^2 + h_1(u)(z-z_0) + h_0(u) = 0.
    \]
    Clearly, solutions are given by
    \[
    z = z_0 - \frac{h_1(u)}2 \pm \sqrt{\frac{h_1(u)^2}4 - h_0(u)}.
    \]
    It remains to study the term $\frac{h_1(u)^2}4 - h_0(u)$. 
    From the relation
    \[
    (u- g(z))^2 - h(z)^2 (z-z_0)^3 = K(z,u)( (z-z_0)^2 + h_1(u)(z-z_0) + h_0(u))
    \]
    it follows (by setting $z=z_0$) that
    $
    h_0(u) = {(u-u_0)^2}/{K(z_0,u)}
    $
    and that (after differentiating with respect to $z$ and setting $z=z_0$)
    \[
    h_1(u) = - \frac{2(u-u_0)g'(z_0) + K_z(z_0,u)h_0(u)}{K(z_0,u)}.
    \]
    Furthermore, by differentiating twice with respect to $z$ and setting
    $z=z_0$ and $u=u_0$ we obtain
    $
    K(z_0,u_0) = g'(z_0)^2.
    $
    This implies that
    \[
    \frac{h_1(u)^2}4 - h_0(u) = O((u-u_0)^3).  
    \]
    By computing few more derivatives we can also compute $K_u(z_0,u_0)$
    which gives 
    \[
    \frac{h_1(u)^2}4 - h_0(u) = C(u-u_0)^3   + O(u-u_0)^4
    \]
    with
    $
    C = \frac{h(z_0)^2}{ g'(z_0)^5} \ne 0.
    $
    This finally implies that
    \[
    \sqrt{\frac{h_1(u)^2}4 - h_0(u)} = H(u) (u-u_0)^{3/2}
    \]
    for an analytic function $H(u)$ with $H(u_0)\ne 0$ and completes the proof of the theorem.
\end{proof}

\end{document}
